\documentclass[11pt]{article}
\usepackage[margin=0.7in]{geometry}
\usepackage{amsmath,amssymb}
\usepackage{enumitem}
\usepackage{tikz}
\usepackage{pgfplots}
\pgfplotsset{compat=1.18}
\usepackage{fancyhdr}
\usepackage{parskip}
\setlist[enumerate]{leftmargin=*, itemsep=6pt}
\pagestyle{fancy}
\fancyhf{}
\fancyfoot[L]{\small\textit{en}Vision\textsuperscript{\textregistered} Algebra 2}
\fancyfoot[C]{\small 2-3 Lesson Quiz, short form}
\fancyfoot[R]{\small\thepage}
\renewcommand{\headrulewidth}{0pt}
\newcommand{\choice}[1]{\tikz[baseline=-0.6ex]\node[draw, circle, inner sep=0.6pt, minimum size=3.3mm, font=\scriptsize\bfseries] {#1};\;}
\newcommand{\checkbox}{$\square$\;}
\newcommand{\workline}{\par\vspace{12pt}\noindent\rule{\linewidth}{0.4pt}\par}

\begin{document}

\noindent\textbf{Name}\; \rule{6cm}{0.4pt} \hfill Period \rule{1.2cm}{0.4pt} \quad Date \rule{2.2cm}{0.4pt}

\vspace{4pt}\noindent\rule{\textwidth}{0.6pt}

\vspace{2pt}\noindent{\LARGE\textbf{2-3 \; Lesson Quiz: Factored Form of a Quadratic Function}}

\noindent{\small short form \hfill Questions keep their Savvas quiz numbers. Show your work.}

\vspace{2pt}\noindent\rule{\textwidth}{0.6pt}

\vspace{6pt}
\textbf{1.} Factor the expression $3x^2-11x-4$.

\vspace{10pt}
\hspace{2em}$\Large(\;\rule{1cm}{0.4pt}\,x\;+\;\rule{1cm}{0.4pt}\;)(\,x\;-\;\rule{1cm}{0.4pt}\;)$

\vspace{4pt}
Work:
\workline
\workline

\vspace{12pt}\noindent\rule{\textwidth}{0.4pt}

\vspace{6pt}
\begin{minipage}[t]{0.42\textwidth}
\textbf{2.} Solve the equation $x^2+7x=30$.

\vspace{6pt}
\begin{itemize}[leftmargin=1.2em, itemsep=6pt]
\item[] \choice{A} $x=15$ and $x=-2$
\item[] \choice{B} $x=-6$ and $x=5$
\item[] \choice{C} $x=-10$ and $x=3$
\item[] \choice{D} $x=10$ and $x=-3$
\end{itemize}
\end{minipage}\hfill
\begin{minipage}[t]{0.54\textwidth}
Show the factored equation and the Zero Product Property step:
\workline
\workline
\workline
\workline
\end{minipage}

\vspace{14pt}\noindent\rule{\textwidth}{0.4pt}

\vspace{6pt}
\textbf{5.} A long-jumper lifts off 3 m after starting his run, and lands 6 m from where he lifted off. When he is 8 m from the start line, he is 5 cm above the ground. Write the equation of a parabola that models his path through the air, where $x$ is his horizontal distance from the start line in m and $y$ is his height, in cm.

\vspace{4pt}
\begin{minipage}[t]{0.42\textwidth}
\begin{itemize}[leftmargin=1.2em, itemsep=6pt]
\item[] \choice{A} $y=-x^2+12x-27$
\item[] \choice{B} $y=-3x^2+4x-9$
\item[] \choice{C} $y=3x^2-4x+9$
\item[] \choice{D} $y=x^2-12x+27$
\end{itemize}
\end{minipage}\hfill
\begin{minipage}[t]{0.54\textwidth}
\raggedright
Zeros: \ $x=\rule{1.2cm}{0.4pt}$ \ and \ $x=\rule{1.2cm}{0.4pt}$

\vspace{4pt}
Point: \ $(\rule{0.9cm}{0.4pt},\,\rule{0.9cm}{0.4pt})$

\vspace{4pt}
Work:
\workline
\workline
\workline
\end{minipage}

\vfill
\noindent\textbf{Scoring:} Question 2 and Question 5 are 10 points each. Question 1 is a bonus of up to 5 points. Total out of 20.

\newpage
\noindent{\LARGE\textbf{2-3 Lesson Quiz, short form: answer key}}

\vspace{4pt}\noindent\rule{\textwidth}{0.6pt}

\begin{enumerate}
\item[\textbf{1.}] \textbf{3, 1, 4: $(3x+1)(x-4)$.} $ac=3(-4)=-12$; the factors $1$ and $-12$ have sum $-11$. $3x^2+x-12x-4=x(3x+1)-4(3x+1)=(3x+1)(x-4)$. Check: $3x^2-12x+x-4=3x^2-11x-4$. Savvas quiz item 1, DOK 2, Skills Review A24, A69. Taught by Example 1, Part B (cut); the packet rehearses the grouping step in Example 3, Part B. \textbf{Bonus 5:} all three blanks correct; 3 for the correct binomials with a sign slip; 0 otherwise.
\item[\textbf{2.}] \textbf{C, $x=-10$ and $x=3$.} $x^2+7x-30=0$, $(x+10)(x-3)=0$, so $x=-10$ or $x=3$. Distractor B solves $x^2+x=30$; D has both signs flipped, the error of reading the constants in $(x+10)(x-3)$ as the zeros. Savvas quiz item 2, DOK 2, A24. Example 3, Part A, and Try It 3a. \textbf{10 points:} 6 for C, 4 for the factored equation and the Zero Product Property step.
\item[\textbf{5.}] \textbf{A, $y=-x^2+12x-27$.} Lift-off at $x=3$; he lands 6 m \emph{after} lift-off, at $x=3+6=9$. So $y=a(x-3)(x-9)$. The point $(8,5)$: $5=a(5)(-1)$, $a=-1$. $y=-(x-3)(x-9)=-x^2+12x-27$. Trap: using $6$ as the second zero. D has $a=+1$ (opens up, below ground between the zeros). Savvas quiz item 5, DOK 2, F50, F58. Example 6 and the packet Practice (Practice \& Problem Solving 45). \textbf{10 points:} 4 for the zeros 3 and 9, 2 for the point $(8,5)$, 4 for A.
\end{enumerate}

\textbf{Items left off this form} (available for a retake or Form B):
\begin{itemize}[itemsep=2pt]
\item Savvas item 3, the projectile \mbox{$h(t)=-16t^2+32t+128$}: answer $4$ seconds ($-16(t-4)(t+2)=0$, and time cannot be negative), DOK 2, A29; it is taught by Example 4 (cut) and its function is the same as Try It 4's, so it cannot be a fresh item for students who did Try It 4.
\item Savvas item 4, intervals where $f(x)=x^2+12x+27$ is positive: answer $-9$, $-3$, that is $(-\infty,-9)$ and $(-3,\infty)$, DOK 2, A70; taught by Example 5 (cut).
\end{itemize}

\textbf{IXL replacement:} SmartScore 70 or above on UB5, CJC and HPB (the starred skills; UB5 feeds the item 1 bonus) replaces this quiz with 20/20.

\textbf{Verify} every key against the Savvas Realize quiz PDF (aga\_24\_a2\_na\_0203\_lq.pdf) on quiz day; this form is built from the TE transcription (a2\_2-3\_TE.tex, printed page 82A).

\end{document}
